% Shared preamble of the v6 method figure candidates (v6_method_a.tex, v6_method_b.tex) and the v6 method zoom.
% Units: 1 unit = 1 bp at print size; the figure is built at its final width (5.5 in = 396 bp), so nothing is scaled.
% Type: labels 7.5 pt sans (Helvetica clone, as the paper's \sffamily), small labels 7 pt, panel titles 9 pt bold,
% math 8.5 pt in Computer Modern (the paper's math), scripts never under 6 pt (\DeclareMathSizes).
% Palette (sca_colors.tex, shared with every plot): SCA green, near amber, far blue, neutral greys. Four hues.
% Assets: ../v6/method_assets (v6_method_assets.py; provenance.json there). Build: v6_method.py.
\documentclass[border=0pt]{standalone}
\RequirePackage{fix-cm}
\usepackage[T1]{fontenc}
\usepackage{helvet}
\renewcommand{\familydefault}{\sfdefault}
\usepackage{amsmath}
\usepackage{graphicx}
\usepackage{tikz}
\usetikzlibrary{arrows.meta, calc, positioning}
\DeclareMathSizes{6.5}{6.5}{6.2}{6.2}
\DeclareMathSizes{7}{7}{6.2}{6.2}
\DeclareMathSizes{7.5}{7.5}{6.2}{6.2}
\DeclareMathSizes{8}{8}{6.2}{6.2}
\DeclareMathSizes{8.5}{8.5}{6.4}{6.2}
\DeclareMathSizes{9}{9}{6.6}{6.2}
\DeclareMathSizes{9.5}{9.5}{7}{6.2}
\DeclareMathSizes{10}{10}{7}{6.2}
% squared sizes in sans text: the superscript printed at 6.4 pt (\textsuperscript would give 5 to 6 pt)
\newcommand{\sq}{\textsuperscript{\fontsize{6.4}{7}\selectfont 2}}
\newcommand{\sqt}{\textsuperscript{\fontsize{6.8}{7}\selectfont 2}}
% Written by scafig.py; the same palette as the matplotlib figures.
\definecolor{scaInk}{HTML}{1F2328}
\definecolor{scaMid}{HTML}{5B6570}
\definecolor{scaRule}{HTML}{C3C9CF}
\definecolor{scaGrid}{HTML}{E6E9EC}
\definecolor{scaGlobal}{HTML}{2F6DB3}
\definecolor{scaGlobalLight}{HTML}{DCE7F3}
\definecolor{scaLocal}{HTML}{E08A2B}
\definecolor{scaLocalLight}{HTML}{F9E6D0}
\definecolor{scaAccent}{HTML}{13865F}
\definecolor{scaAccentLight}{HTML}{D7EEE4}
\definecolor{scaNeutral}{HTML}{F4F1EA}
\definecolor{scaZoom}{HTML}{FFC21A}
\definecolor{scaLoss}{HTML}{B85C00}
\definecolor{scaArea}{HTML}{7A4EA3}
\definecolor{scaAreaLight}{HTML}{E9E0F1}
\definecolor{scaFull}{HTML}{9AA3AC}
\definecolor{scaCrop}{HTML}{E4572E}

\input{method_assets/data.tex}
\definecolor{scaLocalText}{HTML}{A85F12}
\definecolor{scaGlobalText}{HTML}{245A96}
\definecolor{scaAccentText}{HTML}{0E6B4C}
\definecolor{scaPanel}{HTML}{F6F7F8}

\newcommand{\ttl}{\fontsize{9}{10.5}\selectfont\bfseries}
\newcommand{\lab}{\fontsize{7.5}{9}\selectfont}
\newcommand{\sml}{\fontsize{7}{8.4}\selectfont}
\newcommand{\mth}{\fontsize{8.5}{10}\selectfont}
\newcommand{\mtl}{\fontsize{9.5}{11}\selectfont}
\newcommand{\mts}{\fontsize{7.5}{9}\selectfont}

\tikzset{
  every node/.style={inner sep=0pt, outer sep=0pt, text=scaInk},
  arr/.style={-{Stealth[length=3.4pt, width=3pt]}, line width=0.6pt, draw=scaInk},
  garr/.style={arr, draw=scaAccent},
  blk/.style={rounded corners=2pt, line width=0.8pt},
  sblk/.style={blk, draw=scaAccent, fill=scaAccentLight},
  nblk/.style={blk, draw=scaFull, fill=white},
  frame/.style={line width=0.5pt, draw=scaFull},
  axis/.style={line width=0.5pt, draw=scaMid},
  wedge/.style={fill=scaGrid},
}

% ---------------------------------------------------------------------------------------------------- image tiles
% Fractions in data.tex are measured from the top-left corner; \fp{x}{y}{S} maps them into a tile of side S whose
% lower-left corner is the scope origin.
\newcommand{\fp}[3]{({(#1)*(#3)},{(#3)-(#2)*(#3)})}
% Query marker: the query token, drawn as a 2.6 bp black square with a white rim at its fractional position.
\newcommand{\QueryMark}[1]{%
  \pgfmathsetmacro{\qmx}{(\QX+0.5*\QS)*#1}\pgfmathsetmacro{\qmy}{#1-(\QY+0.5*\QS)*#1}%
  \fill[white] (\qmx-1.7,\qmy-1.7) rectangle ++(3.4,3.4);
  \fill[scaInk] (\qmx-1.25,\qmy-1.25) rectangle ++(2.5,2.5);}
% Rectangles of the windows (top-left fraction x, y; side s)
\newcommand{\FracBox}[4]{\fp{#1}{#2}{#4} rectangle \fp{#1+#3}{#2+#3}{#4}}

% \TgNear{x}{y}{S}: the 4096^2 target with the native 64-token window (outside veiled)
\newcommand{\TgNear}[3]{%
  \begin{scope}[shift={(#1,#2)}]
    \node[anchor=south west] at (0,0) {\includegraphics[width=#3bp,height=#3bp]{method_assets/img4096.jpg}};
    \fill[white, opacity=0.62, even odd rule] (0,0) rectangle (#3,#3) \FracBox{\NatX}{\NatY}{\NatS}{#3};
    \draw[line width=1pt, scaLocal] \FracBox{\NatX}{\NatY}{\NatS}{#3};
    \draw[frame] (0,0) rectangle (#3,#3);
    \QueryMark{#3}
  \end{scope}}
% \TgFar{x}{y}{S}{radius}: the 4096^2 target under the native 64 x 64 grid of block-centre keys (every key drawn)
\newcommand{\TgFar}[4]{%
  \begin{scope}[shift={(#1,#2)}]
    \node[anchor=south west] at (0,0) {\includegraphics[width=#3bp,height=#3bp]{method_assets/img4096.jpg}};
    \fill[white, opacity=0.55] (0,0) rectangle (#3,#3);
    \foreach \i in {0,...,63}{\foreach \j in {0,...,63}{
      \fill[scaGlobal] ({(4*\i+2.5)*#3/256},{#3-(4*\j+2.5)*#3/256}) circle[radius=#4];}}
    \draw[frame] (0,0) rectangle (#3,#3);
    \QueryMark{#3}
  \end{scope}}
% \AtNear{x}{y}{S}: 1024^2, emulated near view at r = 4 (its attention, amber; outside the window veiled)
\newcommand{\AtNear}[3]{%
  \begin{scope}[shift={(#1,#2)}]
    \node[anchor=south west] at (0,0) {\includegraphics[width=#3bp,height=#3bp]{method_assets/b_near.png}};
    \draw[line width=1pt, scaLocal] \FracBox{\WinX}{\WinY}{\WinS}{#3};
    \draw[frame] (0,0) rectangle (#3,#3);
    \QueryMark{#3}
  \end{scope}}
% \AtFar{x}{y}{S}: 1024^2, emulated far view at r = 4: 16 x 16 block-centre keys, dot area = its attention share
\newcommand{\AtFar}[3]{%
  \begin{scope}[shift={(#1,#2)}]
    \node[anchor=south west] at (0,0) {\includegraphics[width=#3bp,height=#3bp]{method_assets/b_far.png}};
    \foreach \i in {0,...,15}{\foreach \j in {0,...,15}{
      \fill[scaGlobal!55] ({(4*\i+2.5)*#3/64},{#3-(4*\j+2.5)*#3/64}) circle[radius=0.32];}}
    \foreach \dx/\dy/\dr in \FarDots{\fill[scaGlobal] \fp{\dx}{\dy}{#3} circle[radius={\dr*#3}];}
    \draw[frame] (0,0) rectangle (#3,#3);
    \QueryMark{#3}
  \end{scope}}
% \AtFull{x}{y}{S}: 1024^2, the model's full attention (amber inside the r = 4 window, blue outside)
\newcommand{\AtFull}[3]{%
  \begin{scope}[shift={(#1,#2)}]
    \node[anchor=south west] at (0,0) {\includegraphics[width=#3bp,height=#3bp]{method_assets/b_full.png}};
    \draw[line width=0.6pt, scaLocal, dash pattern=on 1.6pt off 1.2pt] \FracBox{\WinX}{\WinY}{\WinS}{#3};
    \draw[frame] (0,0) rectangle (#3,#3);
    \QueryMark{#3}
  \end{scope}}

% ---------------------------------------------------------------------------------------------------- the fit
% \Quad{x0}{y0}{w}{h}{which}: J(pi) rescaled to [0, 1] on pi in [0, 1] (data.tex), with its minimiser.
% which: 1 = the illustrated head only, 3 = far, illustrated and near heads.
\newcommand{\Quad}[5]{%
  \begin{scope}[shift={(#1,#2)}]
    \begin{scope}[xscale=#3, yscale=#4/1.14, shift={(0,0.10)}]
      \ifnum#5=3
        \draw[line width=0.8pt, scaGlobal, line join=round] plot coordinates {\JcFar};
        \draw[line width=0.8pt, scaLocal, line join=round] plot coordinates {\JcNear};
      \fi
      \draw[line width=1.0pt, scaInk, line join=round] plot coordinates {\JcMid};
    \end{scope}
    \draw[axis] (0,0) -- (#3,0);
    \draw[axis] (0,0) -- (0,#4);
    \ifnum#5=3
      \foreach \p in {\JpFar,\JpNear}{\fill[scaAccent] ({#3*\p},{#4*0.10/1.14}) circle[radius=1.5];}
    \fi
    \fill[scaAccent] ({#3*\JpMid},{#4*0.10/1.14}) circle[radius=1.8];
    \foreach \v/\t in {0/0, 1/1}{
      \draw[axis] ({#3*\v},0) -- ({#3*\v},-1.6);
      \node[anchor=north, text=scaMid] at ({#3*\v},-2.6) {\sml\t};}
    \node[anchor=south, text=scaAccentText] at ({#3*\JpMid},{#4*0.10/1.14+2.6}) {\mts$\pi^\star$};
    \node[anchor=west, text=scaMid] at (#3+2.5,0) {\mth$\pi$};
    \node[anchor=south east, text=scaMid] at (-2.5,#4-6) {\mth$J$};
  \end{scope}}

% ---------------------------------------------------------------------------------------------------- the table
% \Cards{x}{y}{w}{h}{dx}: the r = 2 card behind, the r = 4 card in front (lower-left corner x, y), one block's 24
% heads (rows) by 13 noise bins (columns, sigma falling left to right); the measured cell is outlined in green.
\newcommand{\Cards}[5]{%
  \begin{scope}[shift={(#1,#2)}]
    \node[anchor=south west] at (#5,#5) {\includegraphics[width=#3bp, height=#4bp]{method_assets/table_r2.png}};
    \draw[frame] (#5,#5) rectangle ++(#3,#4);
    \fill[white] (0,0) rectangle (#3,#4);
    \node[anchor=south west] at (0,0) {\includegraphics[width=#3bp, height=#4bp]{method_assets/table_r4.png}};
    \draw[frame] (0,0) rectangle (#3,#4);
    \draw[line width=1pt, scaAccent] ({\TabCol*#3/\TabCols},{#4-(\TabRow+1)*#4/\TabRows}) rectangle
      ++({#3/\TabCols},{#4/\TabRows});
  \end{scope}}
